%% [migrated] heading absorbed into chapter wrapper: Gauge Fields

\section{Yang Mills \& Gauge Invariance}

\begin{itemize}
\item Classical Yang Mills Lagrangian in continuum, Euclidean space
\end{itemize}
\begin{align}
    \mathcal{L}_{YM} &= \frac{1}{4}\sum_a(F_{\mu \nu}^a)^2 \color{green} + \theta \epsilon_{\mu \nu \rho \sigma} F^{\mu \nu}_a F^{\rho \sigma}_a \nonumber \\
    &= \frac{1}{2} \text{tr}(F_{\mu \nu} F_{\mu \nu}) \color{green} + \theta \epsilon_{\mu \nu \rho \sigma} F^{\mu \nu}_a F^{\rho \sigma}_a, 
\end{align}
where we will for now ignore the $\theta$ term in green and where\footnote{Note that we are using  conventions where
\begin{equation}
\begin{aligned}
& \text{tr}\left[T_{a} T_{b}\right]=\frac{1}{2} \delta_{a b} \\
& {\left[T_{a}, T_{b}\right]=i f_{a b c} T_{c}}.
\end{aligned}
\end{equation}}
\begin{align}
F^{\mu \nu} &= \frac{i}{g}[D^\mu, D^\nu] = F^{\mu \nu}_a T_a = \partial^\mu A^\nu - \partial^\nu A^\mu - i g [A^\mu, A^\nu],\quad  D^\mu = \partial^\mu - i g A^\mu \nonumber \\
F^{\mu \nu}_a &= \partial^\mu A^\nu_a - \partial^\nu A^\mu_a + g f_{abc}A^\mu_b A^\nu_c.
\end{align}
for the gluon degree of freedom
\begin{equation}
A_{\mu}(x)=A_{\mu}^{a}(x) T_{a}
\end{equation}

where \(T_{a}\) are the generators of the Lie algebra \(G\)

\begin{itemize}
  \item Under a local transformation
\end{itemize}

\begin{equation}
A_{\mu}(x) \rightarrow A_{\mu}^{\prime}(x)=\Omega(\theta) A_{\mu}(x) \Omega^{-1}(\theta)-\frac{2 i}{g}\left(\partial_{\mu} \Omega(\theta)\right) \Omega^{-1}(\theta)
\end{equation}

with

\begin{equation}
\Omega(\theta)=\exp \left[-\frac{i}{2} \theta^{a}\left(x\right) T_{a}\right] \in \text { Lie group } G
\end{equation}

the Lagrangian above is invariant and

\begin{equation}
D^{\mu} \rightarrow \Omega D^{\mu} \Omega^{-1}, \qquad F^{\mu \nu} \rightarrow \Omega F^{\mu \nu} \Omega^{-1}
\end{equation}

\section{Parallel transporters (Holonomies)}
\begin{itemize}
\item It will be useful to consider \underline{parallel transporters} that translate the coordinates in the internal symmetry space from one spacetime point to another.

\item The general form of a parallel transporter for a path $\Gamma$ between $x$ and $y$ is
\begin{equation}
    E_\Gamma(x,y) = \mathcal{P}\exp\left[i \int_\Gamma A \cdot \text{d}s\right],
\end{equation}
where $\mathcal{P} \exp$ implies the ``path-ordered exponential''.

\item Note that $E_\Gamma(x,y)$ solves the equation
\begin{equation}
\label{eq:par_transp_eq}
    D_sE_\Gamma(x,y) = 0,
\end{equation}
where the path-covariant derivative along the path $z^\mu(s)$ is
\begin{equation}
    D_s = \frac{\text{d}z^\mu(s)}{\text{d}s} D_\mu = \frac{\text{d}z^\mu}{\text{d}s}\partial_\mu + i \frac{\text{d}z^\mu}{\text{d}s}A_\mu = \frac{\text{d}}{\text{d}s} + i \frac{\text{d}z^\mu}{\text{d}s}A_\mu(z). 
\end{equation}
\item Now the above implies
\begin{equation}
\frac{\text{d} E_{\Gamma}(x, y)}{\text{d} s}=-i \frac{\text{d} z^\mu}{\text{d} s} A_{\mu}(z) E_{\Gamma}(x, y) \equiv-M(s) E_{\Gamma}(x, y),
\end{equation}
which is solved by ($M$ is a matrix)
\begin{align}
    E_\Gamma(x = z_\mu(0), y = z_\mu(t)) &= \sum_{m=0}^\infty (-i)^m \int_0^t \text{d}s_m \int_0^{s_m}\text{d}s_{m-1}\ldots\int_0^{s_2}\text{d}s_1 M(s_m)M(s_{m-1})\ldots M(s_1) \nonumber\\
    &\equiv \mathcal{P}\exp \left\{i \int_0^t \text{d}s \frac{\text{d}z_\mu}{\text{d}s} A_\mu(z) \right\} \nonumber \\
    &= \mathcal{P}\exp\left\{ i \int_\Gamma \text{d}z_\mu A^\mu(z) \right\}.
\end{align}
One can (and should)  check that this satisfies Eq.~\eqref{eq:par_transp_eq}.
\end{itemize}
\begin{figure}
    \centering
    \includegraphics[max width=0.5\textwidth, center]{fig-notes-009.jpg}
    \caption{An illustration of one possible path $z_\mu$.}
    \label{fig:parallel_transporter}
\end{figure}
\begin{itemize}
\item Parallel transporters are elements of the gauge group
\item Under a gauge transformation, parallel transporters transform nicely
\begin{equation}
   \boxed{E_\Gamma(x,y) \underset{\text{gauge transform}}{\longrightarrow} \underset{\in G}{\underbrace{\Omega(x)}} E_\Gamma(x,y) \Omega^{-1}(y)}
\end{equation}
\item On a lattice geometry, parallel transporters between neighboring sites assume particular importance and are called link variables:
\begin{align}
    U_\nu(n) &= E_\Gamma(n_\mu, n_\mu + \hat{\nu}) \quad \text{with} \, \Gamma \, \text{a straight link path} \, z_\mu(t) = n_\mu + t\hat{\nu} \nonumber \\
    &= \mathcal{P} \exp \left\{ i \int_n^{n + \hat{\nu}} \text{d}z_\mu A^\mu(z) \right\}.
\end{align}
Note that
\begin{equation}
    \int_n ^{n+\hat{\nu}} \text{d}z_\mu = a \delta_{\mu \nu},
\end{equation}
so that assuming $a \rightarrow 0$, $\Gamma$ is very short and we can assume $A^\mu(z) = A^\mu(n)$, we find
\begin{equation}
    \boxed{U_\nu(n) = \exp \{i a A^\nu(n) \} }
\end{equation}
 \begin{align}
   U_\nu(n) \rightarrow 1 + i a A_\nu(n) + \mathcal{O}(a^2)
\end{align}
\item Parallel transporters satisfy:
\begin{itemize}
    \item $E_{\cancel{\Gamma}}(x,x) = 1$ where $\cancel{\Gamma}$ is the empty path
    \item $E_{\Gamma_2 \circ \Gamma_1}(x,z) = E_{\Gamma_2}(x,y)E_{\Gamma_1}(y,z)$, where $\Gamma_2 \circ \Gamma_1$ denotes composition of paths
    \item $E_{-\Gamma}(y,x) = \left[E_\Gamma(x,y) \right]^{-1}$, where $-\Gamma$ is the path $\Gamma$ traversed in the reverse sense
\end{itemize}
\item An extended path $\Gamma$ can be constructed from link variables
\begin{equation}
E_\Gamma(x, y)=\includegraphics[width=0.1\textwidth]{fig-notes-061.png}=U_{\mu}(x) U_{\nu}(x+\hat{\mu}) \cdots U_{\mu}(y-\hat{\mu})
\end{equation}
\item Parallel transporters around closed paths are called Wilson loops
\begin{equation}
L_{\Gamma}[U]=\operatorname{tr}\left[\prod_{n, \mu \in \Gamma} U_{\mu}(n)\right]
\end{equation}
which are gauge invariant
\item The entirety of YM theory can be built from these objects
\end{itemize}

\section{Gauge Invariance on a lattice}

\begin{itemize}
\item Consider a set of complex scalar fields $\varphi = (\varphi_1, \ldots \varphi_n)^T$
    \begin{itemize}
    \item The corresponding Lagrangian contains
    \begin{equation}
        \mathcal{L} = \partial_\mu \varphi^\dagger \partial_\mu \varphi = \sum_{i=1}^n \partial_\mu \varphi_i^* \partial_\mu \varphi_i,
    \end{equation}
    \item This theory is invariant under global transformations of $\varphi$ under a continuous Lie group $G$
    \begin{equation}
        \varphi(x) \rightarrow \Omega(\theta) \varphi(x).
    \end{equation}
    \item To make this theory invariant under local $G$ transformations in the continuum,
    \begin{equation}
        \varphi(x) \rightarrow \varphi^\prime(x) = \Omega(\theta(x)) \varphi(x),
    \end{equation}
    with
    \begin{equation}
        \Omega(\theta(x)) = \exp \left[ \frac{i}{2}T_a \theta^a(x) \right],
    \end{equation}
    we promote $\partial_\mu$ to a gauge covariant derivative $D_\mu=\partial_\mu+A_\mu$, introducing a new gauge field.
    \end{itemize}
\item The corresponding lattice theory has a kinetic term
\begin{align}
    \mathcal{L}_\text{latt} &\sim \sum_{\nu = 1}^{d+1}\text{Tr}[(\varphi^\dagger(n + \hat{\nu}) - \varphi^\dagger(n))(\varphi(n + \hat{\nu}) - \varphi(n)) ] \nonumber \\
    &\sim \sum_{\nu = 1}^{d+1}\text{Tr}\left[\varphi^\dagger(n + \hat{\nu})\varphi(n + \hat{\nu}) - \varphi^\dagger(n + \hat{\nu}) \varphi(n) - \varphi^\dagger(n)\varphi(n + \hat{\nu}) + \varphi^\dagger(n)\varphi(n) \right].
\end{align}
If we let $\varphi(n) \rightarrow \varphi^\prime(n) = \Omega(\theta(n))\varphi(n)$, we see $\mathcal{L}$ is clearly not invariant and clearly no local modifications can make $\varphi^\dagger(n + \hat{\nu})\varphi(n)$ be invariant.
\item To solve this problem, we can use the link fields $U_\mu(n)$ associated with the links between every pair of neighboring sites in the lattice. In particular $U_\mu(n)$ is associated with the link $(n, n+\hat \mu)$
    \begin{itemize}
        \item Demanding 
        \begin{equation}
            \quad \boxed{U_{\mu}(n) \rightarrow U_{\mu}^{\prime}(n)=\Omega(\theta(n)) U_{\mu}(n) \Omega^\dagger(\theta(n+\hat{\mu}))},
        \end{equation} 
        we can write a gauge invariant Lagrangian
        \begin{equation}
        \mathcal{L}_{\text {latt}} \sim \sum_{\nu = 1}^{d+1}\text{Tr}\left[\varphi^\dagger(n + \hat{\nu})\varphi(n + \hat{\nu}) - \varphi^\dagger(n + \hat{\nu}) U_{-\nu}(n + \hat{\nu})\varphi(n) - \varphi^\dagger(n)U_\nu(n)\varphi(n + \hat{\nu}) + \varphi^\dagger(n)\varphi(n) \right],
        \end{equation}
        where we use $U_{-\nu}(n) = U_\nu(n - \hat{\nu})^\dagger$ for $\nu\in\{1,\ldots,d+1\}$ to denote a link traversed in the opposite direction, so the introduction of the link variables has made the hopping terms gauge invariant. We can check the consistency of this relation with the transformation:
        \begin{equation}
            U_{-\nu}(n) \rightarrow U^\prime_{-\nu}(n) = [ \Omega(n - \hat{\nu}) U_\nu(n - \hat{\nu}) \Omega^\dagger(n)]^\dagger = \Omega(n) U_{-\nu} \Omega^\dagger(n - \hat{\nu}).
        \end{equation}
        \item The transformation property of the link variables imply that they are elements of the group $G^{DV}$
        \item Just as in the continuum formulation, demanding gauge invariance has required the introduction of new fields
    \end{itemize}
  
\end{itemize}


\section{Wilson gauge action}
\begin{itemize}
  \item Wilson's original formulation of LGT is based on the smallest gauge invariant quantity, the $1 \times 1$ plaquette
\end{itemize}
\begin{equation}
\begin{aligned}
U_{\mu v}(n) & =U_{\mu}(n) U_{\nu}(n+\hat{\mu}) U_{\mu}^{\dagger}(n+\hat{\nu}) U_{\nu}^{\dagger}(n) \\
& =U_{\mu}(n) U_{\nu}(n+\hat{\mu}) U_{-\mu}(n+\hat{\mu}+\hat{\nu}) U_{-\nu}(n+\hat{v})
\end{aligned}
\end{equation}

\begin{figure}
\includegraphics[max width=0.5\linewidth]{fig-notes-010.jpg}
\caption{A depiction of the plaquette $U_{\mu \nu}(n)$.}
\end{figure}

\begin{itemize}
  \item Using this, the Wilson action for a \(S U(N)\) gauge symmetry is
\end{itemize}

\begin{equation}
S_{G}[U]=\frac{\beta}{N} \sum_{n \in \Lambda} \sum_{\mu < \nu} \operatorname{Re} \operatorname{tr}\left(1-U_{\mu \nu}(n)\right)
\end{equation}

where

\begin{equation}
\beta=\frac{2 N}{g^{2}}
\end{equation}

\begin{itemize}
  \item In the classical continuum limit (just $a \rightarrow 0$ geomerically) this gives the continuum action 
  \begin{itemize}
      \item To see this, one needs to use BCH to combine the products in $U_{\mu \nu}(n)$
      \begin{equation}
          U_{\mu \nu}(n) = \exp \left\{ i a A_\mu + i a A_\nu(n + \hat{\mu}) - \frac{a^2}{2}[A_\mu(n), A_\nu(n + \hat{\mu})] + \ldots\right\}
      \end{equation}
      \item Now Taylor expand the fields around $n$
      \begin{align}
          A_\mu(n + \hat{\nu}) &= A_\mu(n) + a \partial_\nu A_\mu(n) + \mathcal{O}(a^2) 
      \end{align}
     \begin{align}
          \Rightarrow U_{\mu \nu}(n) &= \exp\left\{i a^2 [\partial_\mu A_\nu(n) - \partial_\nu A_\mu(n)] + i [A_\mu(n), A_\nu(n)] \right\} + \mathcal{O}(a^2) \nonumber \\
          &= \exp\{i a^2 F_{\mu \nu}(n)\} + \ldots\nonumber \\
          &= 1 + i a^2 F_{\mu \nu}(n) - a^4 (F_{\mu \nu}(n))^2 + \ldots \label{eq:Umunu}
      \end{align}
      and thus
      \begin{equation}
          S_G[U] \underset{a\rightarrow0}{\longrightarrow}\frac{a^4 \beta}{N} \sum_{n \in \Lambda}\sum_{\mu, \nu} \operatorname{tr}[F_{\mu \nu}(n)^2] + \mathcal{O}(a^2)
      \end{equation}
      as ${\rm tr}(F_{\mu\nu})=0$.
  \end{itemize}
  \item It is also interesting to interpret Eq.~\eqref{eq:Umunu} as
  \begin{equation}
      U_{\mu \nu}(n) = \exp\{i a^2 F_{\mu \nu}(n)\} =\exp\{i \int_\Box F_{\mu \nu}(n)ds_{\mu\nu}\}
  \end{equation}
  where $ds_{\mu\nu}$ is the integral over the area of the plaquette. In the Abelian case, this would simply arise from the definition of the plaquette as a line integral around the four links forming the boundary and Stokes' theorem. This suggests the interpretation of the plaquette as the exponential of the (chromo-)magnetic flux through the plaquette, ie the non-Abelian analogue of the Aharanov-Bohm phase factor.
\end{itemize}



\section{YM partition function and correlation function}

\begin{itemize}
  \item The path integral formalism is similar to scalar fields
\end{itemize}

\begin{equation}
\begin{aligned}
& Z=\int \mathcal{D} U \exp \left[-S_{G}[U]\right] \\
& \langle\Omega| \mathcal{O}_{a} \mathcal{O}_{b}|\Omega\rangle=\frac{1}{Z} \int \mathcal{D} U \mathcal{O}_{a}[U] \mathcal{O}_{b}[U] e^{-S_{G}[U]}
\end{aligned}
\end{equation}

where the integration measure is

\begin{equation}
\mathcal{D} U=\prod_{n \in \Lambda} \prod_{\mu=1}^{D} \text{d} U_{\mu}(n)
\end{equation}

\begin{itemize}
  \item Since \(U_{\mu}(n)=e^{ia A_{\mu}(n)}\) is a group valued object, the measure \(\text{d} U_{\mu}(n)\) involves integration over all elements of the continuous compact Lie group \(G\)
  \begin{itemize}
        \item Since \(Z\) and \(S_{G}[U]\) are gauge invariant, then the measure must be invariant under gauge transformations: the Haar measure which is defined by requiring that for \(U, V \in G\)
        \begin{equation}
            \underset{\text{L, R invariance}}{\text{d}U = \text{d}(V U) = \text{d}(UV)}, \underset{\text{linearity}}{\quad \int \text{d}U(\alpha f(U) + \beta g(U)) = \alpha \int \text{d}U f(U) + \beta \int \text{d}U g(U)}, \quad \text{and} \underset{\text{normalization}}{\int \text{d}U = 1}
        \end{equation}
        \item Concretely for $U = U(\omega) = \exp(i \omega_a T_a)$, the metric in $G$ is $\text{d}s^2 = g(\omega)_{mn} \text{d}\omega_n \text{d}\omega_m$, where
        \begin{equation}
            g(w)_{mn} = \operatorname{tr}\left[ \frac{\partial U(\omega)}{\partial \omega_n}\frac{\partial U(\omega)^\dagger}{\partial \omega_m}\right],
        \end{equation}
        and the measure is
        \begin{equation}
            \text{d} U = c \sqrt{\operatorname{det}g(\omega)}\prod_a \text{d}\omega_a \quad \text{with constant $c$ such that $\int \text{d}U = 1$}.
        \end{equation}
        For the case of $SU(N)$, see J Broznan Phys Rev D 38 (1988)1994 for explicit constructions. 
  \end{itemize}
  \item The group integration forces the expectation of gauge variant quantities to vanish (Elitzur's theorem, see Elitzur 1975); that is
  \begin{equation}
      \langle \mathcal{O} \rangle = 0 \quad \forall\mathcal{O} \underset{G}{\rightarrow} \mathcal{O}^\prime \neq \mathcal{O}.
  \end{equation}
  \begin{itemize}
      \item To see this, consider a single group integration over a group element $U \in G$:
      \begin{equation}
          \int \text{d} U  U= \int \text{D}(UV) U = \int \text{d}U U V^\dagger = \int \text{d}(W^\dagger U^\prime) U^\prime V^\dagger = \int \text{d}U^{\prime \prime} W U^{\prime \prime}V^\dagger = W \left[ \int \text{d} U U\right] V^\dagger
      \end{equation}
      which holds for all $W$ and $V^\dagger$, implying $\int \text{d}U U = 0$.
      \item Note that any group integral $\int dg f(g)=0$ unless $f(g)$ has a singlet component
      \item Now consider a gauge variant quantity such as a single gauge link (or a product of gauge links)
      \begin{equation}
        \begin{aligned}
        & \left\langle U_{\mu}(n)\right\rangle=\frac{1}{Z} \int \mathcal{D} U e^{-S_{G}[U]} U_{\mu}(n) \\
        & =\frac{1}{Z} \int D \hat{U} e^{-\hat{S}[\hat{U}]} \prod_{\nu= \pm1}^{ \pm d} \text{d} U_{\nu}(n) e^{-\tilde{S}\left[\hat{U},\left\{U_{\nu}(n)\right\}\right]} U_{\mu}(n),
        \end{aligned}
        \end{equation}
        where $\mathcal{D}\hat{U}$ is the measure for all links not connected to site $n$ and $$S_G[U] = \hat S[\hat{U}] +  \underset{\text{all terms involving links at site $n$}}{\underbrace{\tilde{S}[\hat{U}, \{ U_\nu(n)\}]}}$$

        There are $2(D-1)$ plaquettes directly involving $U_\mu(n)$:
        \begin{equation}
            P_{\mu \nu}(n), P_{\mu \nu}(n - \hat{\nu}) \quad \forall \nu \neq \mu
        \end{equation}
        For
        \begin{equation}
        \begin{gathered}
        P_{\mu \nu}(n)=U_{\mu}(n) U_{\nu}(n+\hat{\mu}) U_{\mu}^{\dagger}(n+\hat{v}) U_{\nu}^{\dagger}(n)
        \end{gathered}
        \end{equation}
\begin{figure}
\includegraphics[width=0.6\linewidth, center]{fig-notes-011.jpg}
\caption{An illustration of the links that we will use in our various changes of variables, as well as their associated plaquettes.}
\end{figure}
We can now use the invariance of the Haar measure to change variables for the particular link:
\begin{equation}
U_{\nu}(n) \rightarrow U_{\mu}(n) U_{\nu}(n) \color{green} \quad \circled{1} 
\end{equation}
\begin{equation}
\begin{aligned}
 \Rightarrow \operatorname{tr}\left(P_{\mu \nu}(n)\right) &\rightarrow \operatorname{tr}\left(U_{\mu}(n) U_{\nu}(n+\hat{\mu}) U_{\mu}^{\dagger}(n+\hat{\nu}) U_{\nu}^{\dagger}(n) U_{\mu}^{\dagger}(n)\right) \\
&=\operatorname{tr}\left[U_{\nu}(n+\hat{\mu}) U_{\mu}^{\dagger}(n+\hat{\nu}) U_{\nu}^{\dagger}(n)\right] \quad \text { independent of } U_{\mu}(n)
\end{aligned}
\end{equation}
which is independent of  $U_{\mu}(n)$.
However, this change also affects

\begin{equation}
\begin{aligned}
\quad P_{\mu v}(n-\hat{\mu})&=U_{\mu}(n-\hat{\mu}) U_{\nu}(n) U_{\mu}^{\dagger}(n-\hat{\mu}+\hat{\nu}) U_{\nu}^{\dagger}(n-\hat{\mu}) \\
&  \stackrel{{\tiny{\color{green} \circled{1}}}}{\rightarrow} U_{\mu}(n-\hat{\mu}) U_{\mu}(n) U_{\nu}(n) U_{\mu}^{\dagger}(n-\hat{\mu}+\hat{\nu}) U_{\nu}^{\dagger}(n-\hat{\mu})
\end{aligned}
\end{equation}

So now change variables for the link
\begin{equation}
  U_{\mu}(n-\hat{\mu}) \rightarrow U_{\mu}(n-\hat{\mu}) U_{\mu}^{\dagger}(n) \quad {\color{green}\circled{2}}  
\end{equation}
So

\begin{equation}
P_{\mu v}(n-\hat{\mu})\stackrel{{\tiny{\color{green} \circled{2}}}}{=} U_{\mu}(n-\hat{\mu}) U_{\nu}(n) U_{\mu}^{\dagger}(n-\hat{\mu}+\hat{\nu}) U_{v}^{\dagger}(n-\hat{\mu})
\end{equation}
independent of \(U_{\mu}(n)\)

Also,
\begin{equation}
P_{\mu\nu}(n-\hat{\nu})=U_{\mu}(n-\hat{\nu}) U_{\nu}(n-\hat{\nu}+\hat{\mu}) U_{\mu}^{\dagger}(n) U_{\nu}^{\dagger}(n-\hat{\nu})
\end{equation}

under \(\left(U_{\nu}^{\dagger}(n-\hat{\nu})\left(=U_{-\nu}(n)\right) \rightarrow U_{\mu}(n) U_{\nu}^{\dagger}(n-\hat{\nu})\right)\) \color{green} \circled{3} \color{black} \(\left[U_{\nu}(n-\hat{\nu}) \rightarrow U_{\nu}(n-\hat{\nu}) U_{\mu}^{\dagger}(n)\right]\)

\begin{equation}
P_{\mu \nu}(n-\hat{\nu}) \rightarrow U_{\mu}(n-\hat{\nu}) U_{\nu}(n-\hat{\nu}+\hat{\mu}) U_{\nu}^{\dagger}(n-\hat{\nu}) 
\end{equation}
indep of $U_{\mu}(n)$, and at the same time
\begin{equation}
\begin{aligned}
 P_{\mu \nu}(n-\hat{\nu}-\hat{\mu})&=U_{\mu}(n-\hat{\nu}-\hat{\mu}) U_{\nu}(n-\hat{\nu}) U_{\mu}^{\dagger}(n-\hat{\mu}) U_{\nu}^{\dagger}(n-\hat{\nu}-\hat{\mu}) \\
& \rightarrow U_{\mu}(n-\hat{\nu}-\hat{\mu}) U_{\nu}(n-\hat{\nu}) U_{\mu}^{\dagger}(n) U_{\mu}(n) U_{\mu}^{\dagger}(n-\hat{\mu}) U_{\nu}^{+}(n-\hat{\nu}-\hat{\mu}) \\
& =U_{\mu}(n-\hat{\nu}-\hat{\mu}) U_{\nu}(n-\hat{\nu}) U_{\mu}^{\dagger}(n-\hat{\mu}) U_{\nu}^{\dagger}(n-\hat{\nu}-\hat{\mu}) \quad \text { indep of } U_{\mu}(n)
\end{aligned}
\end{equation}

The same procedure proceeds in the other \(D-2\quad(\mu, \nu)\) planes and by changing variables on the \(2D-1\) links from site \(n\) other than \(U_{\mu}(n)\) the action has been made independent of \(U_{\mu}(n)\) !! Consequently

\begin{align}
    \langle U_\mu(n)\rangle &= \frac{1}{Z}\int \mathcal{D} \hat{U} e^{- \hat{S}[\hat{U}]} \Pi_{\nu = \pm 1, \nu \neq \mu}^{\pm d} \text{d} U_\nu (n) e^{-\tilde{S}[\hat{U}, \{U_\nu(n)\}_{\nu \neq \mu}]} \int \text{d}U_\mu(n) U_\mu(n) \nonumber \\
    &= 0.
\end{align}
  \end{itemize}


  \item Note that the demonstrated independence of the action on any one link is simply a manifestation of the redundancy of the gauge link degrees of freedom
  \item For a compact Lie group, the integration for each link is over a compact region and so is finite for nonsingular integrands. Therefore using link variables removes the need to fix a gauge. Nevertheless the gauge can be fixed in numerous ways and is necessary if doing lattice regulated perturbation theory from a lattice gauge theory action.
  \begin{itemize}
      \item Landau gauge $(\partial_\mu A_\mu^a = 0)$: for smooth fields this is equivalent to minimizing $\int \text{d}x \operatorname{tr}[A_\mu(x)^2]$ which is achieved on link fields by minimizing the functional (for $SU(N)$)
      \begin{equation}
          L_U[\Omega] = \frac{1}{N D V}\sum_{\mu, n} \operatorname{Re}\operatorname{tr} [\Omega(n) U_\mu(n) \Omega^\dagger(n + \hat{\mu})],
      \end{equation}
      and for Coulomb gauge the sum is just over spatial $\mu$. To see this, let $\{U_\mu\}$ be the configuration that maximizes $L_U[\Omega]$ on a given gauge orbit. Then for $\Omega = 1 + i \omega(n)$ near the identity,
      \begin{equation}
          L_U[1 + i \omega(n)] \simeq L_U[1] + \frac{1}{2 N D V} \sum_{\mu, n} i w^a(n) \operatorname{Tr}[T^a(U_\mu(n) - U_\mu^\dagger(n)) - (U_\mu(n - \hat{\mu}) - U_\mu^\dagger(n - \hat{\mu}))T^a] 
      \end{equation}
      By definition \(U_{\mu}\) is a stationary point so
        \begin{equation}
            0 = \frac{\partial L}{\partial w^a(n)} = \frac{i}{2NDV} \sum_\mu \operatorname{Tr}[T^a(U_\mu(n) - U_\mu^\dagger(n)) - (U_\mu(n - \hat{\mu}) - U_\mu^\dagger(n - \hat{\mu}))T_a] 
        \end{equation}
      and noting that $U_\mu(n) - U_\mu^\dagger(n) \approx 2 i a A_\mu\left(n + \frac{\hat{\mu}}{2}\right) + \mathcal{O}(a^3)$
      \begin{align}
          \Rightarrow 0 &= \sum_\mu \operatorname{Tr}\left[\left(A_\mu\left(n + \frac{\hat{\mu}}{2}\right) - A_\mu\left(n - \frac{\hat{\mu}}{2}\right)\right)T^a\right] \nonumber \\
          &= \sum_\mu \partial_\mu A_\mu^a + \mathcal{O}(a)
      \end{align}
      Note that finding $\text{max}_\Omega L_U[\Omega]$ is a challenging optimization task over $(N^2 - 1)V$ degrees of freedom
      \item even finding a local maximum is computationally demanding - see Olivera \& Silva, Comp Phys Com 158 (2004)73 for discussion.
      \item since $\partial_\mu A^\mu = 0$ intersects the gauge orbit of a given gauge field multiple times (Gribov copies), there are multiple maxima of $L_U[\Omega]$.
      \item For gauge invariant quantities, the presence of Gribov copies (Gribov-Singer, Singer proved you cannot define a gauge fixing that avoids the ambiguity on $S^4$) does not matter as they are included in the integral defining expectation values and in $Z$. For gauge variant quantities, it is possible to define a subsidiary condition to choose a global maximum (e.g. minimizing $\int \text{d}x \operatorname{tr} [A_\mu(x)^2]$) but it isn't clear if this means the same as in the continuum.
      
  \end{itemize}
  \item Temporal gauge $(A_4(x) = 0)$ is useful, for example in showing reflection positivity of Wilson gauge action. In terms of link variables temporal gauge corresponds to $U_4(n) = 1$.

  This is a specific case of a maximal tree gauge, in which any set of links $\{ U_\mu(n)\}_{\text{tree}}$ that contains no closed loops can be set to unity (or to any fixed group elements).

  As we saw in calculating $\langle U_\mu(n) \rangle = 0$, the action depends on plaquettes rather than links and gauge freedom can be used to remove dependence on given links.

  The simplest (?) gauge transformation acts only on a single site and modifies all $2D$ links that touch it and in general can be chosen to set any one of those links to unity (see Fig.~\ref{fig:simplest_gauge_transformation}.
\end{itemize}
\begin{figure}
\includegraphics[width=\linewidth, center]{fig-notes-012.jpg}
\caption{A demonstration of the gauge transformation on one site which could be used to set one of the gauge links to zero.
\label{fig:simplest_gauge_transformation}
}
\end{figure}
It is clear that a gauge transform at a given site can set at most one link to unity, but by implementing particular transformations at each site a given gauge field $\{ U_\mu(n)\}$ can be transformed to $\{ U_\mu (n)\}_\text{tree}$ where all links in a tree (a graph on $\Lambda$ with no closed loops) are unity; it is easy to see that one cannot gauge fix away closed loops, as the trace of a plaquette is gauge invariant and so cannot be set to unity by a gauge transformation. An example of a maximal tree gauge is shown in Fig.~\ref{fig:maximal_tree_gauge}.

\begin{figure}
\includegraphics[width=0.6\linewidth, center]{fig-notes-014.jpg}
\caption{A maximal tree gauge; note that any any additional red link would create a cycle.
\label{fig:maximal_tree_gauge}
}
\end{figure}

Temporal gauge fixing sets all temporal links to unity,
\begin{figure}
\includegraphics[width=0.5\textwidth, center]{fig-notes-013.jpg}
\caption{An illustration of temporal gauge in $D=2$.}
\end{figure}
Note that temporal gauge fixing is not a maximal tree gauge.

The above has assumed a lattice of infinite extent in each dimension.
\begin{figure}
\includegraphics[width=0.5\textwidth, center]{fig-notes-015.jpg}
\caption{Temporal gauge with periodic boundary conditions, $U_{-4}^{\dagger}\left(n_{t}=0\right)=U_{4}\left(n_{t}=N_{t}-1\right)$.
}
\end{figure}


On a finite lattice, fixing all temporal lints to unity would form a closed loop and cannot be done because 
\begin{equation}
P(\vec{n})=\prod_{t} U_{4}(\vec{n}, t)
\end{equation}
is gauge invariant 

But one can set

\begin{equation}
\begin{aligned}
& U_{4}(\vec{n}, t)=1 \mid 0 \leqslant t<N_{t}-1 \\
& U_{4}\left(\vec{n}, N_{t}-1\right)=P(\vec{n}).
\end{aligned}
\end{equation}

An interesting fact: $N_{max tree}\sim e^{zV}$, where for $D=4$ it is known that $z \sim 2.00$ (see Shrock and Wu, cond-mat/0004341 ).

\section{Observables in Y M theory}

\begin{itemize}
  \item Gauge invariant quantities whose values can be determined in YM theory include
  \begin{itemize}
      \item Wilson loops: closed paths of gauge links
      \begin{equation}
          \langle W (n; \underset{\mu_i \in \{ \pm1, \ldots, \pm D\}}{\underbrace{\mu_1, \mu_2, \ldots, \mu_n}})\rangle = \langle \operatorname{Tr}[U_{\mu_1}(n)U_{\mu_2}(n+\hat{\mu}_1)\ldots U_{\mu_n}(n - \hat{\mu}_n)] \rangle
      \end{equation}
      \item Polyakov loops: straight paths in the $\hat{\mu}$ direction that close around boundary
      \begin{equation}
          \mathcal{P}_\mu(\vec{n}_\perp) = \left\langle \operatorname{Tr} \left[ \prod_{i=0}^{N_\mu - 1} U_\mu(\vec{n}_\perp, n_\mu = i)\right] \right\rangle \qquad \text{where $\vec{n}_\perp$are coords transverse to $\hat{\mu}$},
      \end{equation}
      \begin{equation}
          \mathcal{P}_\mu = \sum_{\vec{n}_\perp} \mathcal{P}_\mu(\vec{n}_\perp)
      \end{equation}
      \item Correlations between Wilson loops or Polyakov loops, e.g.
      \begin{equation}
    \sum_{\vec{n}, \vec{m}}\left\langle W\left(\vec{n}, 0 ; \mu_{1}, \ldots, \mu_{n}\right) W\left(\vec{m}, \hat{t}; \nu_{1}, \ldots, \nu_{p}\right)\right\rangle
    \end{equation}
    which can be used to determine the masses of excitations in pure YM: glueballs
    \item In $SU(n > 2)$ can also use Levi-Civita to make observables, as shown in Fig.~\ref{fig:lc_observables}
    \begin{equation}
    \varepsilon_{a b c} \varepsilon_{a^\prime b^\prime c^\prime}\left[W\left(x, \mu_{1}, \ldots, \mu_{j}\right)\right]_{a a^{\prime}}\left[\tilde{W}\left(x, \nu_{1}, \ldots, \nu_{k}\right)\right]_{b b^{\prime}}\left[\tilde{\tilde{W}}\left(x, \rho_{1}, \ldots ,\rho_{l}\right)\right]_{c c^{\prime}}
    \end{equation}
    \begin{figure}
        \centering
        \includegraphics[width=0.5\linewidth]{fig-notes-060.png}
        \caption{A schematic depiction of an observable created using the Levi-Civita symbol}
        \label{fig:lc_observables}
    \end{figure}
  \end{itemize}
\end{itemize}
\begin{figure}
\includegraphics[width=0.5\linewidth, center]{fig-notes-062.png}
\caption{An illustration of the correlation function between two Wilson loops at different times.}
\end{figure}



\section{RGE and the continuum limit}

\begin{itemize}
  \item Consider any physical quantity with the dimensions of mass \(H=H(g(a), a)\) written as a function of the coupling and cutoff [Here we consider pure YM]
  \begin{itemize}
  \item We demand cutoff independence
  \begin{equation}
  \label{eq:cutoff_ind}
  0=a \frac{d}{d a} H(g(a), a)=a \frac{\partial H}{\partial a}+a \frac{\partial g}{\partial a} \frac{\partial H}{\partial g}
  \end{equation}
  \item Since $[H] = 1$, $H = a^{-1}f(g(a))$ for dimensionless $f$
  \begin{equation}
      \Rightarrow a \frac{\partial H}{\partial a} = a(- a^{-2}f) = - H.
  \end{equation}
  \item Substituting in Eq.~\eqref{eq:cutoff_ind},
  \begin{equation}
      \beta_{CS} \equiv a \frac{\partial g}{\partial a} = \frac{\partial g}{\partial \log a} = \frac{H(g(a),a)}{\partial H/\partial g}
  \end{equation}
  \item We know from continuum PT that
  \begin{equation}
      \beta_{CS}(g) = - \beta_0 g^3 - \beta_1 g^5
  \end{equation}
  with (here we specify to $SU(N_c)$ and allow for $N_f$ flavours of quarks)

  \begin{equation}
      \beta_0 = \frac{1}{16 \pi^2}\left[ \frac{11 N_c}{3} - \frac{2 N_f}{3}\right], \quad \beta_1 = \left(\frac{1}{16 \pi^2}\right)^2 \left[\frac{34}{3}N_c^2 - \frac{10}{3}N_c N_f - \frac{N_f}{N_c}(N_c^2 - 1) \right]
  \end{equation}
  which are universal (valid with a lattice regulator) - higher orders are not and depend on the choice of $H$ and the regulator
  \end{itemize}
  \item The RGE can be solved
  \begin{align}
      &\frac{\partial \log a}{\partial g} = \frac{1}{-\beta_{CS}(g)} \Rightarrow \int_{g_0}^{g(a)} \frac{\text{d}\log a}{\text{g}^\prime}\text{d}g^\prime = \int_{g_0}^{g(a)}\frac{1}{-\beta_{CS}(g^\prime)}\text{d}g^\prime \nonumber \\
      &\Rightarrow \log a + \log \tilde{\Lambda}_L = \int_{g_0}^{g(a)}\frac{\text{d}g^\prime}{\beta_0 g^{\prime 3}} = \frac{1}{2 \beta_0}\left(\frac{1}{g^2(a)}-\frac{1}{g_0^2}\right) \nonumber \\
      &\Rightarrow a(g) = \frac{1}{\Lambda_L}\exp(-\frac{1}{2 \beta_0 g^2})\qquad \text{where $\Lambda_L$, $\tilde{\Lambda}_L=a(g_0)$ are integration constants.}
  \end{align}
  Keeping the $g^5$ term yields $a(g) = \frac{1}{\Lambda_L}\exp\left(-\frac{1}{2 \beta_0 g^2}\right) (\beta_0 g^2)^{-\beta_1/2\beta_0^2}$.
  \item Inverting the RGE gives
  \begin{equation}
      \frac{1}{g(a)^2} = - 2 \log(a\Lambda_L) + \frac{\beta_1}{\beta_0}\log(\log(a^{-2}\Lambda_L^{-2}))+\ldots
  \end{equation}
  so $a \rightarrow 0$ requires $g(a) \rightarrow 0$ (and $\beta \rightarrow \infty$) so the continuum limit is the limit of vanishing bare coupling with physical quantity $H$ held fixed
\end{itemize}

\section{Transfer matrix and Kogut-Susskind Hamiltonian}

\begin{itemize}
  \item The QM interpretation of LGT is provided by defining the Hilbert space and the transfer matrix / Hamiltonian.
  \begin{itemize}
      \item Note that $S_G[U]$ is reflection positive so a Hermitian Hamiltonian and Hilbert space exist
      \item A basis for the Hilbert space can be constructed from eigenstates of operators corresponding to spatial link variables
      \begin{equation}
          \hat{U}_m(\vec{n})|U\rangle = U_m(\vec{n})|U\rangle \qquad m = 1,\ldots, d, \quad \vec{n}\in\Lambda_d
      \end{equation}
      Note that this is a matrix equation in the group $G$. If links are coordinatised by a set of real parameters $\alpha_k(\vec{n})$ the above is equivalent to labeling the bases states $|\alpha\rangle$ such that $\hat{\alpha}_k(\vec{n})|\alpha\rangle = \alpha_k(\vec{n}) |\alpha\rangle$.
      \item These states are orthonormal and complete
      \begin{align}
          \langle U^\prime|U\rangle = \delta(U - U^\prime) &= \prod_{\vec{n}, m}\delta(U_m(\vec{n}),U^\prime_m(\vec{n})) \nonumber \\
          1 &= \int\mathcal{D}U |U\rangle \langle U |.
      \end{align}
      \item Our goal is to express the partition function as
      \begin{equation}
      Z = \int \mathcal{D}U e^{S_G[U]} = \operatorname{Tr}(\hat{T}^{N_t})
      \end{equation}
      Now
      \begin{equation}
          \operatorname{Tr}(\hat{T}^{N_t}) = \prod_{t=0}^{N_t - 1}\left[\int \prod_{\vec{n}, m}\text{d}U_m(\vec{n},t)\right]\langle U_{t+1}| \hat{T} | U_t \rangle,
      \end{equation}
      where the states are states in the timeslice denoted by the respective subscript.
      \item To find $\hat{T}$ , we need to break the action up into various pieces involving or not temporal links. Note that
      \begin{align}
          S_G[U] &= \# \sum_n \sum_{\mu<\nu} \operatorname{Re}\operatorname{Tr}P_{\mu \nu}(n)\nonumber \\
          &= \frac{2}{g^2}\left( \frac{a}{a_t}\sum_n \sum_{m=1}^d \operatorname{Re}\operatorname{Tr} P_{m4}(n) + \frac{a_t}{a}\sum_n \sum_{m=1}^d\sum_{l < m}\operatorname{Re}\operatorname{Tr}P_{lm}(n)\right)
      \end{align}
      where the first term comes from time-space plane plaquettes, the second term comes from plaquettes in spatial planes, and we make the lattice spacing in the $t$ direction different. We could also write
      \begin{equation}
          = \frac{\beta}{N}\{\xi\sum_{P_t}\operatorname{Re}\operatorname{Tr}(P_t) + \frac{1}{\xi} \sum_{P_s}\operatorname{Re}\operatorname{Tr}(P_s) \}
      \end{equation}
      where $\xi=a/a_t$ is the anisotropy and the sums are over the space-time plaquettes and teh space-space plaquettes, respectively..
      \item Comparison with the above shows
      \begin{equation}
          \hat{T} = \exp\left(-\frac{a_t}{2}\hat{W}\right) \hat{T}_K^\prime \exp\left(-\frac{a_t}{2}\hat{W}\right),
      \end{equation}
      where
      \begin{equation}
          \hat{W} = \frac{-2}{g^2 a}\sum_{P_s}\operatorname{Re}\operatorname{Tr}(P_s),
      \end{equation}
      and
      \begin{equation}
      % [migrated] duplicate label eq:99 removed (eq:transf_gauge kept)
          \langle U_{t+1} | \hat{T}_K^\prime | U_t \rangle = \prod_{\vec{n}} \prod_{m=1}^d \int \text{d} U_4(\vec{n},t) \exp{\left[\frac{2a}{g^2 a_t}\operatorname{Re}\operatorname{Tr} \{ U_m(\vec{n},t)\underset{\widetilde{U}^\dagger_m(\vec{n},t+1)}{\underbrace{U_4(\vec{n}+\hat{m},t) U^\dagger_m(\vec{n},t+1) U^\dagger_4(\vec{n},t)}}\} \right]}
          \label{eq:transf_gauge}
      \end{equation}
      \begin{figure}
          \centering
          \includegraphics[width=0.5\linewidth]{fig-notes-063.png}
          \caption{An illustration corresponding to the exponent in Eq.~\eqref{eq:transf_gauge}.}
          \label{fig:enter-label}
      \end{figure}
      Now, the term $U_4(\vec{n}+\hat{m}, t)U_m^\dagger(\vec{n},t+1)U_4^\dagger(\vec{n},t)$ can be viewed as a gauge transformation of all the spatial links on timeslice $t+1$ by $\Omega(\vec{n}) = U_4(\vec{n}, t)$ and all such time-independent gauge transforms are integrated over by $\int\prod_{\vec{n}}\text{d}U_4(\vec{n},t)$ in Eq.~\eqref{eq:transf_gauge}.

      Thus if $|U_{t+1}\rangle$ corresponded to a state that was gauge variant at site $\vec{n}$ or $\vec{n} + \hat{m}$ then the integral over the gauge group will vanish
      \begin{equation}
          \Rightarrow \text{only gauge invariant states exist in the Hilbert space}\nonumber
      \end{equation}
      Projection to the gauge invariant space of physical states can be expressed as 
      \begin{equation}
          \hat{P}_0 = \int \prod_{\vec{n}}\text{d}\Omega(\vec{n}) \hat{D}(\Omega),
      \end{equation}
      where $\hat{D}(\Omega)|U\rangle = |U^{\Omega^\dagger}\rangle$ is the gauge transformed state (and $\hat{D}(\Omega)$ is the operator that applies the gauge transform $\Omega$). The operators $\hat{D}(\Omega)$ provide a unitary representation of the group of time independent gauge transforms on Hilbert space. Clearly $\hat{P}_0$ commutes with the potential $\hat{W}$ above as that involves purely spatial links. 

      The kinetic operator can be written as $\hat{T}_K^\prime = \hat{T}_K \hat{P}_0 (=\hat{P}_0 \hat{T}_K)$ with
      \begin{equation}
          \langle U_{t+1}| \hat{T}_K | U_t \rangle = \prod_{\vec{n}} \prod_{m=1}^d \exp \left[\frac{2a}{g^2 a_t} \operatorname{Re}\operatorname{Tr}\{ U_m(\vec{n}, t)U^\dagger_m(\vec{n}, t+1)\} \right].
          \label{eq:TkME}
      \end{equation}
      Note that this is the form you would immediately get in temporal gauge. $\hat{T}_K$ is a positive operator and defines a Hermitian kinetic operator $\hat{T}_K = \exp\{-a_t \hat{K}\}$ and so the full transfer operator can be written as
      \begin{equation}
          \hat{T} = \hat{P}_0 e^{- a_t \hat{H}} = e^{-a_t \hat{H}}\hat{P}_0.
      \end{equation}
  \end{itemize}
  \item To construct the Hamiltonian we consider the unitary operators that implement changes to a particular link $(\vec{n},j)$ on a timeslice
  \begin{equation}
      R_{\vec{n} j}(\Omega)|U\rangle=\mid U^{\prime} \rangle_{\vec{n} {j}}
  \end{equation}
  such that
  \begin{equation}
      \hat{U}_k(\vec{m})|U^\prime\rangle_{\vec{n} \jmath} = \begin{cases}U_{k}(\vec{m})\left|U^{\prime}\right\rangle_{\vec{n}, \jmath} & \text { elsewhere } \\ \left.\Omega U_{k}(\vec{m}) | U^{\prime}\right\rangle_{\vec{n} \jmath} & \text { for } k=\jmath \text { \& } \vec{m}=\vec{n}\end{cases}
  \end{equation}
  i.e. a transformation on the link $U_\jmath(\vec{n})$ only, where $|U\rangle$ is a field eigenstate on a timeslice.
    Note the above equation is written in terms of the abstract group action $\Omega\circ U_k(\vec{m})$ but more concretely the link $U_m(\vec{n})\to D_\Omega(\vec n) U_m(\vec{n}) D_\Omega^\dagger(\vec n+ \hat m)$ where the $D_\Omega(\vec n)$ are matrices in the fundamental representations of $G$. Note also that $R_{\vec{n} \jmath}(\Omega)$ only alters the $(\vec n,j)$ link.

  
  \begin{itemize}
      \item Clearly
      \begin{equation}
       R_{\vec{n} \jmath}(\Omega) R_{\vec{n} \jmath}\left(\Omega^{\prime}\right)=R_{\vec{n} \jmath}\left(\Omega \Omega^{\prime}\right)
       \end{equation}
       and so these operators furnish a representation of $G$ which is unitary because  the invariance of the measure (See Creutz PRD 15(1977) 1128 for a reference) guarantees that the inner product is unchanged
       \begin{align*}
           \langle \phi | \psi \rangle &= \int \mathcal{D}U\langle \phi | U \rangle \langle U |\psi \rangle = \int \mathcal{D}U \phi^*(U) \psi(U)\\
           &= \int \mathcal{D}(VU)\langle \phi | U \rangle \langle U | \psi \rangle \\
           &= \int \mathcal{D} U \langle \phi | V^\dagger U \rangle \langle V^\dagger U | \psi \rangle = \int \mathcal{D}U \phi^*(V^\dagger U)\psi(V^\dagger U) \\
           &= \int \mathcal{D} U \phi_V^*(U) \psi_V(U) \\
           &= \langle \phi_V | \psi_V \rangle
       \end{align*}
       \item The $\hat{R}_{\vec{n}j}$ ``translate'' the variable $U_\jmath(\vec{n})$ in group space and is generated by the canonically conjugate momenta to the links which will be introduced below.
       \item Now the full $\hat{T}$ can be written as
       \begin{equation}
           \hat{T} = \left[ \prod_{\vec{n} \jmath} \int \text{d}\Omega \hat{R}_{\vec{n}\jmath}(\Omega) \exp \left[\frac{2a}{g^2 a_t}\operatorname{Re}\operatorname{Tr}[\Omega] \right]\right] \exp \left[\frac{2 a_t}{g^2 a}\sum_{P_s}\operatorname{Re}\operatorname{Tr}(P_s) \right]\hat{P}_0
       \end{equation}
       To see this, take the matrix element of one of the terms in the first bracket
       \begin{align*}
           \langle U^\prime|\int\text{d}\Omega \hat{R}_{\vec{n}\jmath}(\Omega) \exp[\alpha \operatorname{Re}\operatorname{Tr}[\Omega]] |U\rangle
           &= \int \text{d}\Omega \exp[\alpha \operatorname{Re}\operatorname{Tr}[\Omega]]\langle U^\prime | \hat{D}_{\vec{n} \jmath}(\Omega)|U\rangle \\
           &= \int \text{d}\Omega \exp[\alpha \operatorname{Re}\operatorname{Tr}[\Omega]]\langle U^\prime | \Omega U \rangle \\
           &= \exp[\alpha \operatorname{Re}\operatorname{Tr}[U U^{\dagger \prime}]]
       \end{align*}
       which is the form in Eq.~\eqref{eq:TkME} (note we used that $\langle U^\prime | \Omega U \rangle =\delta(\Omega-U^\prime U^\dagger)$).
       \item We can parameterize the fundamental representation of group elements $\Omega$ in terms of $\Omega = \exp(i \omega_a T_a)$ and write the Haar measure in terms of the $\omega_a$
       \begin{equation}
           \text{d}\Omega = J(\omega) \prod_a \text{d}\omega_a
       \end{equation}
       where the Jacobian can be written down for a given group $G$; $J(\omega)$ must be regular near $\omega_a = 0$ and $J(\omega) = J(-\omega)$.
       \item Since the $\hat{R}_{\vec{n}j}$ give a representation of $G$, they can be written in terms of generators for that representation
       \begin{equation}
           \hat{R}_{\vec{n}j}(\Omega) = \exp(i \omega_a \hat{l}_{\vec{n}j}^a)
       \end{equation}
       where since $\hat{R}_{\vec{n}j}$ are translations in the $U_{\vec{n}j}$ variables, $\hat l_{\vec{n}j}$ are canonical momentum operators.

       They satisfy
       \begin{align}
           [\hat l_{\vec{n}j}^a, \hat l_{\vec{m}k}^b] &= i f^{abc}\hat l_{\vec{n}j}^c \delta_{\vec{n},\vec{m}}\delta_{j k} \\
           [\hat l_{\vec{n}j}^a, \hat U_{\vec{m}k}] &= - T^a \hat U_{\vec{n}j} \delta_{\vec{n}, \vec{m}}\delta_{j k}\\
           [\hat l_{\vec{n}j}^a, \hat U_{\vec{m}+\hat{k},-k}] &=\hat U_{\vec{m}+\hat{k},-k}T^a \delta_{\vec{n},\vec{m}} \delta_{j k}\\
           [\hat l^2_{\vec{n}j}, \hat l^a_{\vec{m}k}] &= [\hat l^2_{\vec{n}j}, \hat R_{\vec{m}k}] = 0 
       \end{align}
       where $\hat l^2_{\vec{n}j} = \sum_a \hat l_{\vec{n} j}^a$ is quadratic Casimir
       and  the Casimir doesn't depend on link direction, i.e.
       \begin{equation}
           \hat l^2_{\vec{n}j} = \hat l^2_{\vec{n}+\hat{j},-j}
       \end{equation}
       The operators for different links commute with each other, and the $\hat l_{\vec{n}j}^a$ are differential operators on the group parameters.
       \item Now using this
       \begin{equation}
        \hat{T}=\left[\prod_{\vec{n}j} \int \prod_{a} d \omega^{a} J(\omega) e^{\left(\omega^{b} \hat l^b_{\vec{n} j}\right)}  \exp \left\{\frac{2 a}{g^2 a_{t}} \operatorname{Tr}\left(\cos \left(\omega^{b} T^{b}\right)\right)\right\}\right] \exp \left[\frac{2 a_{t}}{g^{2} a} \sum_{P_{s}} \operatorname{Re} \operatorname{Tr}\left(P_{s}\right)\right] \hat{P}_0
        \end{equation}
        and as $a_t \rightarrow 0$ this is dominated by $|\omega|$ near $0$ so as to maximize
        \begin{equation}
            \operatorname{Tr} \cos (\omega T) = N - \frac{1}{4}\omega^2 + \mathcal{O}(\omega^4),
        \end{equation}
        where this holds for any $SU(N)$.
        \item The group integrals above are Gaussian in the $\omega \rightarrow 0$ region and easily integrated to give
        \begin{equation}
        \hat{T}=X \exp \left[-a_t\left\{\frac{g^{2}}{2 a} \sum_{\vec{n},j} \hat{l}_{\vec{n}j}^{2}+\frac{2}{g^{2} a} \sum_{P_{S}} \operatorname{Re} \operatorname{Tr}\left[\hat{P}_{S}\right]\right\}\right] \hat{P}_{0}
        \end{equation}
        where $X$ is an irrelevant constant and $\hat{H}_{KS}$, the contents of the $\{ \}$, is defined to be the Kogut-Susskind Hamiltonian.
        \item The first term contains the canonical momenta and taking $a \rightarrow 0$
        \begin{align*}
            \hat{H}_{KS} &= \sum_{\vec{n}} \left( \frac{g^2}{2} \sum_j \hat{\Pi}_j^2 + \frac{1}{4 g^2} \sum_{l,m = 1}^d \hat{F}_{lm}^2\right)\\
            &= \sum_m\left(\frac{1}{2}\hat{E}^2 + \frac{1}{2}\hat{B}^2\right)
        \end{align*}
        with $E_i = F_{0i} = g \frac{\partial}{\partial t}A_i$ and $B_i = \frac{1}{g} \epsilon_{ijk}F_{jk}$ since $\Pi_j = \frac{\delta L}{\delta \dot{A}_j} = \frac{1}{g^2} F_{0j}$.
        \item The temporal links in $S$ gave us the condition that physical states are gauge invariant under time-independent gauge transformations but there is still redundancy to such gauge transformations in $\hat{H}_{KS}$; alternatively, temporal gauge is not a complete gauge fixing.

        The operator to perform a time-indep gauge transformation  at site $\vec{n}$ is
        \begin{equation}
            \hat J_{\vec{n}}(\Omega) = \prod_j \hat{R}_{\vec{n}\jmath}(\Omega),
        \end{equation}
        where the product is over all $\pm$ directions around site $\vec{n}$.

        All physical states should be invariant under this operation
        \begin{align*}
            \hat J_{\vec{n}}(\Omega)|\psi_{\text{phys}}\rangle &= |\psi_{\text{phys}}\rangle \quad \forall \vec{n}, \Omega\\
            \Rightarrow \sum_j &\hat l_{\vec{n} j}^a |\psi_{\text{phys}}\rangle = 0
        \end{align*}
        which says that the sum of each color of electric fluxes into a site $\vec{n}$ vanishes. This is a discrete rendition of Gauss's Law.
        \item Gauss's Law in temporal gauge removes one d.o.f. per group generator leaving $A_\mu^a$ with $D-2$ degrees of freedom, i.e. $2$ d.o.f. in $D=4$.
  \end{itemize}
\end{itemize}



\section{Strong coupling expansion}

\begin{itemize}
  \item An interesting tool in analysing lattice discretisations of QFTs is the strong coupling expansion in which observables are expanded in powers of $\beta = \frac{2N}{g_0^2}$
  \item This is also called a high temperature expansion and used a lot in statistical systems
  \item Consider the rectangular Wilson loop \includegraphics[width=0.2\linewidth]{fig-notes-064.png}
\begin{equation}
    W_{\cal C} = W(\vec{n}, \vec{r}, t) = \operatorname{tr}\left[S(\vec{n}, \vec{r}, 0) T(\vec{n} + \vec{r}, t) S^\dagger(\vec{n}, \vec{r}, t)T^\dagger(\vec{n}, t) \right]
\end{equation}
with
\begin{equation}
    S(\vec{n}, r \hat{j}, t) = \prod_{\vec{m} = \vec{n}}^{\vec{n} + r \hat{j}} U_j(\vec{m}, t), \quad T(\vec{n}, t) = \prod_{j=0}^t U_4(\vec{n}, j)
\end{equation}
for simplicity.
    \item Expectation value
\begin{align}
    \langle W_\mathcal{C} \rangle &= \frac{1}{Z}\int \mathcal{D} U \exp \left[- \frac{\beta}{N} \sum_{P} \operatorname{Re}\operatorname{Tr}(1-P) \right] \operatorname{tr}\left[ \prod_{l \in \mathcal{C}} U_l\right]\nonumber \\
    &= \frac{1}{Z^\prime} \int \mathcal{D}U \exp \left[- \frac{\beta}{2N} \sum_{P} (\operatorname{tr}(P) + \operatorname{tr}(P^\dagger)) \right] \operatorname{tr}\left[ \prod_{l \in \mathcal{C}} U_l\right],
\end{align}
where $\mathcal{C}$ is the contour constituting the rectangular Wilson loop and we dropped the constant 1 term in the trace in the integral and simultaneously rescaled $Z$ to $Z^\prime$.

    \item If $\beta = 0$ then $\langle W_\mathcal{C} \rangle_{\beta = 0} = \frac{1}{Z^\prime} \int \mathcal{D} U \operatorname{tr}[\prod_{l \in \mathcal{G}}U_l] = 0$ since many links have $\int dU U =0$.

    $\Rightarrow$ In the strong coupling limit, Wilson loops generically vanish which is just a statement that in this limit, sites decouple from each other and behave independently
    \item Expanding the exponential
\begin{equation}
\exp [\ldots]=\sum_{i, j=0}^{\infty} \frac{1}{i!j!}\left(\frac{\beta}{2 N}\right)^{i+j}\left(\sum_{P}\operatorname{tr} P\right)^{i}\left(\sum_{P}\operatorname{tr} P^{\dagger}\right)^{j}.
\end{equation}

Note that for $SU(2)$ there is a factor of $1/2$ since plaquettes are not oriented and $\int dU U_{ab} U_{cd} \neq 0$.

    \item Now in $Z^\prime$, the leading term is from $i=j=0$ and the next from $i=j=1$ (not $i=1, j=0$ as $\int \text{d}U U = 0$).
\begin{equation}
    Z^\prime = 1 + \mathcal{O}(\beta^2)
\end{equation}

    \item For nonzero $\langle W_C \rangle$ we must have $U$ and $U^\dagger$ associated with any link in order $\int dU \ldots \neq 0$, so we need to tile the loop with plaquettes

\includegraphics[max width=\textwidth, center]{fig-notes-016.jpg}

\begin{equation}
    \langle W_C \rangle = \frac{1}{Z^\prime}\int \mathcal{D}U \frac{1}{N_A !} \left( \frac{\beta}{2 N}\right)^{N_A} \left ( \sum_{P \in A_\mathcal{C}} \operatorname{tr} [ P^\dagger]\right)^{N_A}\operatorname{tr}\left[\prod_{l \in \mathcal{C}}U_l\right]
\end{equation}
where $P \in A_\mathcal{C}$ are all plaquettes inside contour $\mathcal{C}$ (i.e. all plaquettes in the loop) and there are $N_A!$ possible arrangements of the $\operatorname{tr}[P^\dagger]$s
\begin{equation}
    \Rightarrow \langle W_\mathcal{C} \rangle = \frac{1}{Z'}\int \mathcal{D} U \left( \frac{\beta}{2 N} \right)^{N_A} \prod_{P \in A_\mathcal{C}}\operatorname{tr}[P^\dagger]\operatorname{tr}\left[\prod_{l \in \mathcal{C}} U_l\right]
\end{equation}

Now we need to integrate each link. We do it in the following order using
\begin{equation}
    \int \text{d}U_m(\hat{n}, t)(U_m(\hat{n}, t))_{ab}(U_m^\dagger(\hat{n},t))_{cd} = \frac{1}{N}\delta_{ad}\delta_{bc}
\end{equation}
\includegraphics[max width=\textwidth, center]{fig-notes-058.png}\\
\begin{align}
\langle W_\mathcal{C} \rangle &= \operatorname{tr}[1] \left( \frac{\beta}{2N}\right)^{N_A}\left(\frac{1}{N}\right)^{N_A}\nonumber \\
&= N \exp \left[ N_A \log \frac{\beta}{2 N^2} \right]
\end{align}

    \item For a loop of size $R \times T$
\begin{align}
    \langle W_{R \times T} \rangle &\sim \exp\left[ R T \log \frac{\beta}{2 N^2}\right] + \mathcal{O}(\beta^2) \nonumber \\
    &\sim \exp[-\sigma A]
\end{align}
where $A = R T a^2$ is the loop area and $\sigma = - \frac{1}{a^2} \log \frac{\beta}{2 N^2}$ is a positive quantity

That is, Wilson loops decrease with an area law!
\begin{figure}
    \centering
    \includegraphics[width=0.5\linewidth]{fig-notes-065.png}
    \caption{A large Wilson loop defining the static quark potential.}
    \label{fig:sqp}
\end{figure}
\item Large Wilson loops give a definition of the static quark potential, as depicted in Fig.~\ref{fig:sqp}

    The loop corresponds to the gauge invariant creation of a static source in the anti-fundamental rep at $(\vec x=0, t=0)$ and a fundamental rep at $(\vec R,0)$ and their propagation in time a distance $T$

    The energy of that system as a function of separation $R=|\vec{R}|$ is the static quark potential
    
\begin{equation}
    V(R) = - \lim_{T\to \infty} \frac{1}{T}\log \langle W_{R \times T} \rangle
\end{equation}
so for $SU(N)$ we have that $V(R) \underset{\beta \to 0}{\rightarrow} \sigma R$, where $\sigma$ is the string tension.

Thus $V(R)$ is the energy cost of separating a quark and antiquark a distance $R$ is linear in the distance, which is linear confinement.


\item We can compute higher orders in the strong coupling expansion, but this is made more efficient through the consideration of group characters which let us deal  with the case of Wilson lines in more complicated representations of the gauge group (and Wilson loops that wind multiply though a given link).

\end{itemize}



\section{An interlude on group characters}

\begin{itemize}
\item As a prelude, some aspects of representation theory are needed.
\item Given a compact Lie group $G$, the representations of group elements $\Omega \in G$ are matrices $D_r(\Omega)$ where $r$ labels the representation
\item The \emph{character} of representation $r$ is the function $\chi_r(\Omega) \equiv \operatorname{Tr}[D_r(\Omega)]$ and its value on the group identity $d_r = \chi_r(1)$ is the dimension of the representation
\item Characters are \emph{class functions} of the group (i.e. they take the same value for all group transformations of a given element)
\end{itemize}

\begin{equation}
\chi_{r}\left(\bar{\Omega} \Omega \bar{\Omega}^{-1}\right)=\chi_{r}(\Omega) \quad \forall \Omega, \bar{\Omega} \in G
\end{equation}

\begin{itemize}
  \item If $\mathcal{D}_G$ is the set of all inequivalent unitary irreducible representations of $G$, then for $D_r, D_s \in \mathcal{D}_G$
\begin{equation}
    \int \text{d}\Omega \chi^*_r(\Omega) \chi_s(\Omega) = \delta_{rs},
\end{equation}
i.e. the characters are orthonormal.

    \item Characters are also complete in the sense that
\begin{equation}
    \sum_{D_r \in \mathcal{D}_G} \chi^*_r(\Omega) \chi_r(\Omega^\prime) = \delta(\Omega^\prime \Omega^{-1})
\end{equation}
where the delta function satisfies $\int \text{d} \Omega^\prime f(\Omega^\prime) \delta(\Omega^\prime \Omega^{-1}) = f(\Omega)$ and $\delta(\Omega^{-1}) = \delta(\Omega)$.
 
    \item \emph{Peter-Weyl theorem}: the $\chi_r$ form an orthonormal basis for all square-integrable class functions on $G$, therefore one can decompose a class function as
\begin{equation}
f(\Omega)=\sum_{D_{r} \in \mathcal{D}_G} f_{r} \chi_{r}(\Omega)
\end{equation}
with
\begin{equation}
f_{r}=\int \text{d} \Omega \chi_{r}^{*}(\Omega) f(\Omega)
\end{equation}
as coefficients; i.e. we can do Fourier analysis on a group manifold.

\item Characters for representations of semi-simple Lie groups are tabulated.
    \item For $G = SU(2)$, we know the irreps can be labeled by $j = 0, 1/2, 1, 3/2, \ldots$ with $d_j = 2j + 1$
\begin{itemize}
    \item We can parameterise $SU(2)$ group elements in terms of generators as 

\begin{equation}
\Omega=\exp (i \varphi \vec{n} \cdot \vec{T})=\cos \varphi+i \vec{n} \cdot \vec{T} \sin \varphi, \qquad 0 \leqslant \varphi \leqslant 2 \pi, \qquad|\vec{n}|=1
\end{equation}

where \(\vec{T}\) are the generators of the corresponding Lie algebra

    \item The corresponding measure
\begin{equation}
\text{d} \Omega=\frac{1}{4 \pi^{2}} \text{d} \varphi \sin ^{2} \frac{\varphi}{2} \text{d}S(\vec{n}),
\end{equation}
where $\text{d}S(\vec{n})$ is the uniform measure on $S^2$ and $\int d S(\vec{n})=4 \pi$.

  \item By explicit calculation of
\begin{equation}
\begin{aligned}
\Omega^{\prime} \Omega \Omega^{\prime-1} & =(\cos \rho+i \vec{m} \cdot T \sin \rho)(\cos \varphi+i \vec{n} \cdot \vec{T} \sin \varphi)(\cos \rho-i \vec{m} \cdot T \sin \rho) \\
& =\cos \varphi-i \vec{n} \cdot T \sin \varphi-2 i \cos \rho \sin \rho \sin \varphi(\vec{m} \times \vec{n}) \cdot \vec{T}
\end{aligned}
\end{equation}
it is clear that class functions must be \(f(\Omega)=f(\varphi)\) ( ie only a function of $\varphi$ with no \(\vec{n}\) dependence in this parameterisation).

  \item Without proof the characters are 
\begin{equation}
    \chi_r(\Omega) = \sin((j + 1/2) \varphi) / \sin(\varphi/2),
\end{equation}
which is the Weyl character formula, and for a class function $f(\Omega)$, the coefficients $f_r$ are given by
\begin{align}
    f_r &= \int \frac{1}{4\pi^2}\text{d}\varphi \sin^2\frac{\varphi}{2}\text{d}S(\vec{n}) \frac{\sin((j+1/2)\varphi)}{\sin(\varphi/2)}f(\varphi) \nonumber \\
    &= \frac{1}{\pi} \int_{0}^{2 \pi} \text{d} \varphi \sin \frac{\varphi}{2}\sin((j+1/2)\varphi)f(\varphi)
\end{align}
\end{itemize}

\end{itemize}


\section{Back to strong coupling}

\begin{itemize}
  \item Beyond this leading order, \(\beta \rightarrow 0\) calculations get more involved. Some corrections to consider are
\begin{enumerate}
    \item Shift one tiled plaquette out of the plane and add four plaquettes from action $(\sim \beta^4)$
    \\ \includegraphics[width=0.3\linewidth]{fig-notes-066.png}
    \item Add two additional plaquettes to the tilings either where there are existing plaquettes or elsewhere $(\sim \beta^2)$ (since $\int dU U_{ab} U_{cd} U^\dagger_{ef} U^\dagger_{gh}=\#(\delta\delta\ldots)\ne0$)
    \\ \includegraphics[width=0.3\linewidth]{fig-notes-067.png}
    
    \item Reverse orientation of one plaquette and add $N-2$ additional plaquettes with same orientations $(\sim \beta^N)$
    (since $\int dU U_{a_1b_1} U_{a_2b_2}\ldots U_{a_Nb_N}=\frac{1}{N!}\varepsilon_{a_1\ldots a_N}\varepsilon_{b_1\ldots b_N}\ne0$)
    \\ \includegraphics[width=0.3\linewidth]{fig-notes-068.png}
    \item Terms coming from expansion of $\frac{1}{Z} = 1 - \mathcal{O}(\beta^2)$ \\ \includegraphics[width=0.3\linewidth]{fig-notes-069.png}
\end{enumerate}
    \item One aspect of tracking all of these is the issue of multiple plaquettes at the same point can be handled with the character expansion

\begin{equation}
\exp [-S[U]]=\exp \left[-\frac{\beta}{2 N} \sum_{\square} \operatorname{Re} \operatorname{Tr}[\square]\right]=\prod_{\square} e^{-\frac{\beta}{2 N} \operatorname{Re} \operatorname{Tr}[\square]}
\end{equation}

Focusing on a single plaquette, \(e^{\left(-\frac{\beta}{2 N}(\square + \square^\dagger)\right)}\) is a class function, so

\begin{equation}
\begin{aligned}
 e^{-S(\square)}&=\exp \left(-\frac{\beta}{2 N} \operatorname{Re} \operatorname{Tr} \square\right)=\sum_{r} f_{r} X_{r}(\square) \\
& =N(\beta)\left(1+\sum_{r \neq 0} \tilde{f}_{r} \chi_{r}(\square) \right) \equiv N(\beta)\left(1+f(\square) \right) \\
\end{aligned}
\end{equation}
(note that $\square$ without trace is in the adjoint+singlet representation) where
\begin{align}
    N(\beta) &= \int \text{d}U e^{-\frac{\beta}{2 N}  \operatorname{Re}\operatorname{Tr}(U)} \nonumber\\
    \tilde{f}_r &= \frac{1}{N(\beta)} \int \text{d}U e^{-\frac{\beta}{2N}\operatorname{Re}\operatorname{Tr}(U)} \chi^*_r(U),
\end{align}
and where $r \neq 0$ means that we do not sum over the singlet representation.

  \item In terms of this
\begin{equation}
\begin{aligned}
\exp \left(\frac{-\beta}{2 N} \sum_{\square} \operatorname{Re}\operatorname{Tr}\square\right) & =\prod_{\square} e^{-S(\square)}=\prod_{\square}(1+f(\square)) \quad \text { dropping } N(\beta) \text { which will cancel in  } \frac{1}{Z} \\
& =\sum_{\mathcal{S}} \prod_{\square \in \mathcal{S}} f(\square),
\end{aligned}
\end{equation}
where $\mathcal{S}$ is a set of plaquettes and the sum is over all such sets (including the empty set).

Each set can be decomposed into connected components (plaquettes are connected if they share a link (but not if they only share a corner), $\mathcal{S} = \bigcup_i X_i$ where $X_i$ and $X_j$ are disconnected if $i \neq j$.

Considering integrals over $e^{-S}$ such as
\begin{equation}
Z=\int \mathcal{D}U e^{-S[U]}=\sum_{\mathcal{S}} \int \mathcal{D}U\prod_{\square \in \mathcal{S}} f(\square),
\end{equation}
the integrations over the $X_i \in \mathcal{S}$ factorize and 
\begin{equation}
    Z = \sum_{\mathcal{S} = \sum_i X_i}\prod_i \Phi(X_i), \qquad \Phi(X_i) = \int \mathcal{D}U \prod_{\square \in X_i} f (\square),
\end{equation}
where only the active links in $X_i$ contribute to $\Phi(X_i)$ (see M{\"u}nster Nucl Phys B 180(1981) 23).

  \item Since each $\square$ comes with a $\beta/2N$, the larger a set $\mathcal{S}$ is, the more suppressed it will be: $\beta^{|\mathcal{S}|}$ but there are a rapidly growing $\#$ of large sets so the series is likely asymptotic
  \item Exploring the cluster expansion is an exercise! Note that apart from the values of $f(\square)$, the expansion is independent of the gauge group.
  
  \item In $SU(2)$, $r = 0,1/2, 1, 3/2, \ldots$ and $\chi_r$ is given before

\begin{equation}
\begin{aligned}
\tilde{f}_{r} & =\frac{1}{2 \pi} \int_{0}^{2 \pi} d \varphi \sin ^{2} \frac{\varphi}{2} 2 \exp \left(\beta \cos \frac{\varphi}{2}\right) \frac{\sin \left(\left(r+\frac{1}{2}\right) \varphi\right)}{\sin \varphi / 2} \\
& =2\left(2{r+1}\right) I_{2 r+1}(\beta) / \beta
\end{aligned}
\end{equation}
where the  modified Bessel function  $I_{n}(z) \sim \frac{1}{\Gamma(n+1)}\left(\frac{z}{2}\right)^{n}\left(1+\theta\left(z^{2}\right)\right)$
\begin{equation}
\Rightarrow f(\square)  =\sum_{r \neq 0} \frac{2}{\beta}(2 r+1) I_{2 r+1}(\beta) \frac{\sin \left(\left(r+\frac{1}{2}\right) \varphi\right)}{\sin \varphi / 2}
\end{equation}

However we want to expand in $\beta$ so don't keep all representations.

In $SU(2)$ the leading cluster correction is the type $1$ above which adds four plaquettes relative to the tiled loop and so is $\mathcal{O}(\beta^4)$ so adding this contribution
\begin{equation}
    \langle W_{R\times T} \rangle = 2 \exp(-\sigma_0 A) \{ 1 + 4 A \left(\frac{\beta}{4}\right)^4 + \ldots\},
\end{equation}
where $4A$ is the number of bumps which contribute.

Graphs that count with higher powers of $A$ come in but it has been shown that they exponentiate:

\begin{equation}
\left\langle W_{R \times T}\right\rangle=2 \exp (-\sigma A+\mu(R+T))
\end{equation}

Correction have been computed up to \(\beta^{14}\) [Münster] with many types of decoration of diagrams needed

\begin{equation}
\sigma=\ln \frac{\beta}{4}+4\left(\frac{\beta}{4}\right)^{4}+\frac{176}{3}\left(\frac{\beta}{4}\right)^{8}+\cdots
\end{equation}

  \item In \(SU(3)\) the representations are \(r=1,3, \overline{3}, 6, \overline{6}, 8,10, \overline{10}, \ldots\)
  or alternatively labeled by Dynkin coefficients $(p,q) = (1,1), (2,1), (1,2), (3,1), (1,3), (2,2), (4,1), (1,4), \ldots$ with dimensions 
\begin{equation}
    d_r = \frac{1}{2}p q(p+q)
\end{equation}

and for

\begin{equation}
\begin{aligned}
& U=R\left(e^{i A}, e^{i B}, e^{-i(A+B)}\right) R^{\dagger} \in D_{(2,1)}(G)
\end{aligned}
\end{equation}
(note class functions only depend on $A$ and $B$, so we can do a rotation to $U$ to diagonalize it and dont care what the rotation matrix $R$ is)
\begin{equation}
    \chi_{p,q}(U) = -\frac{i(e^{i p A - i q B} - e^{- i q A + i p B} + e^{-ip(A+B)}(e^{-iqA} - e^{-iqB}) + e^{i q (A+B)}(e^{i p B} - e^{i p A}))}{\underset{s(A,B)}{\underbrace{8 \sin \frac{A-B}{2} \sin \frac{A+2B}{2} \sin \frac{2A+B}{2}}}}
\end{equation}

e.g.
\begin{equation}
    \chi_{2,1} = e^{iA} + e^{iB} + e^{-i(A+B)}
\end{equation}
And the measure is

\begin{equation}
d U=s(A, B) d A d B \quad-\pi \leqslant A, B \leqslant \pi \quad \text { for } U \in SU(3)
\end{equation}

So the orthonormality condition on characters is

\begin{equation}
\int \frac{d A d B}{(2 \pi)^{2}} s^{2}(A, B) \chi_{p, q}^{*}(A, B) \chi_{p^\prime, q^{\prime}}(A, B)=\delta_{p p^{\prime}} \delta_{q q^\prime}
\end{equation}

This leads to very demanding calculations but in the end the results are similar

\begin{equation}
\sigma=\ln \frac{\beta}{18}+4\left(\frac{\beta}{18}\right)^{4}+12\left(\frac{\beta}{18}\right)^{5}+\ldots
\end{equation}
\begin{figure}
    \centering
    \includegraphics[width=0.5\linewidth]{fig-notes-070.png}
    \caption{A depiction of screening, showing a tube of fundamental plaquettes along the perimeter of the large loop.}
    \label{fig:screen}
\end{figure}
\item Some general notes are in order
\begin{enumerate}
    \item The same analysis works for $\mathcal{G} = U(1), Z_2, \ldots$ and in particular linear confinement is seen in the strong coupling limit of these theories
    \item If we construct a Wilson loop using links in the adjoint representation of $SU(N)$ with components $\mathfrak{U}_{ab} = 2 \operatorname{tr}(U^\dagger T_a U T_b)$, this has perimeter law scaling since tiling the perimeter with a tube of fundamental plaquettes yields $\int \text{d} U U_{ab}U^\dagger_{cd}\mathfrak{U}_{kl} = \frac{2}{N^2 - 1}(T_k)_{ac}(T_l)_{bd} \neq 0$, which means that the potential between adjoint sources is asympotically constant at $R \to \infty$; such sources are screened by gluons. Specifically, it is irreps where $p - q \cong 0$ mod $N$ that can be screened, as depicted in Fig.~\ref{fig:screen}
    
    \item Similar expressions can be (and have been) performed for other quantities such as 
    \begin{equation}
    \langle W(r \times s, t=0) W^*(r \times s, t = \tau)\rangle \underset{\tau \to \infty}{\to} \# \exp(-E_{\text{glueball}} \tau)
    \end{equation}
    which allows strong coupling calculation of glueball masses
    
    \item The strong coupling expansion breaks down at some point so the above is not a proof of confinement in QFT! How it breaks down depends on the gauge group
    \begin{itemize}
        \item $U(1)$: at $\beta \sim 1$ there is a phase transition ($1^{st}$ order, see phase diagram in Fig.~\Ref{fig:strong_coupling_u1}) from a confining phase to a Coulomb phase where $V(r) \sim \frac{1}{r} + \text{const}$. Alan Guth proved this has to happen (using a different action) but calculations show where it happens.
        \item For $SU(2)$ and $SU(3)$ numerical evidence shows no Coulomb phase at any $\beta$, but the strong coupling expansion breaks down at some $\beta_r$ (see phase diagram Fig.~\Ref{fig:strong_coupling_su23}). The transition at $\beta_r$ is infinite order and known as the roughening transition and is indicative of discretisation being a dominant effect.
        \item For $SU(N>3)$ there is a phase transition
        \item All of this depends on the particular lattice action!
    \end{itemize}
\end{enumerate}
\begin{figure}
    \centering
    \includegraphics[width=0.5\linewidth]{fig-notes-059.png}
    \caption{The phase diagram for the $U(1)$ theory showing where the strong coupling expansion breaks down due to a first order phase transition.
    \label{fig:strong_coupling_u1}}
\end{figure}
\begin{figure}
    \centering
    \includegraphics[width=0.5\linewidth]{fig-notes-095.pdf}
    \caption{The phase diagram for the $SU(2)$ and $SU(3)$ theories showing where the strong coupling expansion breaks down due to the roughening transition.
    \label{fig:strong_coupling_su23}}
\end{figure}
\end{itemize}

\section{Gauge equivariant quantities}

L-CNNs

\section{Wilson flow}

TODO
